\section{Early Warning Systems : Current Landscape and Issues}
Machine learning has led to state-of-the-art early warning alert systems for many high-stakes applications, from public health to finance to earthquake safety \citep{escobar2020automated, mousavi2019cred, al2020application}.
In this thesis, we are motivated by applications in healthcare, especially recent early warning alert systems for critical care hospital settings ~\citep{hylandEarlyPredictionCirculatory2020,sendakRealWorldIntegrationSepsis2020,wellnerPredictingUnplannedTransfers2017}.
In such settings, alerts are intended to trigger extra attention from clinical staff on specific patients predicted to be at risk of adverse outcomes.
Clinicians have limited time that could have many productive uses; it is critical that automated alerts identify patients who truly need help and do not suggest those who do not need attention.

An unneeded alert, also known as a \emph{false alarm} or a \emph{false positive}, has two detrimental consequences. In the moment, it pulls resources away from where they could be better used. In the long term, too many false alarms can cause clinical staff to distrust the alerts all together, a phenomenon known as \emph{alarm fatigue}~\citep{cvach2012monitor,deb2015alarm,sendelbach2013alarm}.
Some false alarms are inevitable, especially in clinical tasks where the adverse outcome to predict is quite rare.
In the mortality risk task we study later, only about 8-10\% of patients die in the hospital, and most survive many hours before death.
If a system is able to issue alerts throughout a patient's stay, it is critical that the model is designed to avoid alarm fatigue and deliver overall net benefit.
If the tool fails to limit false alarms to an acceptable rate, alerts may be ignored completely.

The key challenge of designing alert systems is thus to \emph{balance} the false alarm rate (related to \emph{precision}) with the true positive rate (known as \emph{recall})~\citep{romero-brufauWhyCstatisticNot2015}.
In the clinic, recall measures the fraction of truly at-risk patients correctly identified by an alert.
Unfortunately, standard objectives for training early-warning classifiers, such as binary cross entropy (BCE), are not designed specifically to address false alarms.
Many early warning systems for clinical settings \citep{hylandEarlyPredictionCirculatory2020,futomaLearningDetectSepsis2017} are trained using such objectives and only balance false alarm concerns in a secondary threshold selection or early stopping stage \emph{after} training.
Such post-hoc adjustment is limited and may not identify the ideal tradeoff between recall and precision.
Our experiments show the deficiency of binary cross entropy even with post-hoc adjustment or per-class weighting.

To overcome this challenge, we show that any binary classifier trainable via stochastic gradient descent (SGD) can, by changing its loss function, be steered toward solutions that offer better control of false alarms. 
Our proposed objective seeks to maximize the number of truly at-risk patients who are helped by alerts (maximize recall), subject to meeting a guarantee on the fraction of all alarms that will be false.
%Our work builds on previous technical work by \citet{ebanScalableLearningNonDecomposable2017} that  proposed a tractable way to maximize recall subject to a precision constraint.
%Our work offers tighter bounds that substantially improve performance and better optimization algorithms. We further offer rigorous empirical validation of these bounds on realistic clinical tasks for the first time.

% often designed for settings with \emph{balanced} ratios of positive and negative examples.
%While this is widely understood bAs is widely understood and further demonstrated in our toy data experiments,

In this chapter, we present two key contributions:
 
%that directly improve interpretability of machine learning in healthcare:
%\begin{enumerate}%[leftmargin=0cm,itemindent=.4cm,labelwidth=\itemindent,labelsep=0cm,align=left]

\textbf{1. New tractable bounds for maximizing recall subject to a minimum precision constraint.} 
While it is easy to compute a model's precision on a concrete dataset, it is not tractable to compute gradients with respect to precision, making optimization difficult.
Previous work by~\citet{ebanScalableLearningNonDecomposable2017} suggested surrogate bounds based on the hinge loss that are amenable to SGD.
However, we find these bounds are too loose and lead to suboptimal performance. 
We thus derive a family of bounds based on the logistic sigmoid function that can be made arbitrarily tight.
We find our new bounds achieve far better performance.


\textbf{2. Demonstration of empirical success on a realistic clinical task.} 
While our constrained optimization objective has been suggested before for general-purpose settings~\citep{ebanScalableLearningNonDecomposable2017}, we find that for clinical early warning systems it is particularly suitable yet not thoroughly studied.
Recent methods for false alarm control \citep{fathonyAPperfIncorporatingGeneric2020,tsoiHeavisideFunctionApproximation2021} focus evaluation on small, repurposed tabular datasets, not realistic scenarios for early warning system deployment.
This leaves open the critical question of performance in real-world clinical settings.
Our experiments on two large EHR datasets suggest that our proposed objective can lead to absolute gains of 0.01-0.15 in recall while satisfying the desired precision constraint(see Sec.~\ref{sec:mimic_mortality_prediction} and \ref{sec:mimic_mortality_prediction_deployment}). 
%Our chosen objective is further \emph{interpretable} to stakeholders in the target application, who can specify a maximum false alarm rate that would be tolerable in their daily practice.

We view these contributions as critical steps to achieving effective early warning systems in clinical settings. We emphasize that none of our technical innovations are specific to healthcare; early warning systems in any domain can benefit from our approach.